% /root/transport; GPT-6 Astra Pro; 2026-11-22 UTC.
% Effort inherited; numerical setting not exposed.
\section{HD-F1-T07: the original canonical norm surface has a smooth Abel image}
\label{sec:f1-norm-abel-embedding}

Fix a connected etale cyclic cover
\[
 \pi:Z\longrightarrow C,\qquad \operatorname{Gal}(Z/C)=\langle\sigma\rangle
       \simeq\mathbb Z/3,
\]
where $C$ is a smooth nonhyperelliptic curve of genus three.
Then $g(Z)=7$. Put
\[
 Q=|K_C|\simeq\mathbb P^2,\qquad
 T=(\pi^{(4)})^{-1}Q\subset\operatorname{Sym}^4Z,\qquad
 a:\operatorname{Sym}^4Z\longrightarrow\operatorname{Pic}^4Z.
 \tag{F1T7.1}
\]
The degree convention on a Prym torsor is the restricted
Jacobian polarization $\beta$ of type $(1,1,3,3)$.

The following conclusions hold for these actual maps. The surface
$T$ is the disjoint union of exactly three smooth connected
surfaces $T_j$. Each $T_j\to Q$ has degree $27$. The reduced
Abel image
\[
                     S_j=a(T_j)
 \tag{F1T7.2}
\]
is a smooth embedded surface in the corresponding connected norm
fibre, a torsor under the four-dimensional Prym.
The map $T_j\to S_j$ is the blowup of finitely many distinct
points, and its exceptional curves are exactly its complete
degree-four pencils. In particular its generic degree is one.
The exact polarization degree is
\[
                       \int_{S_j}\beta^2=36.
 \tag{F1T7.3}
\]
An exact count of these exceptional curves is not required for
these conclusions. The conormal rank on $S_j$ is a further
calculation; the statement here supplies the specified smooth
embedding for that calculation.

\paragraph{1. Smoothness and the actual conormal sequence.}
We recall and independently check the differential calculation
of the fixed-surface contribution. Write
$F=\operatorname{Nm}_\pi\circ a$. Then $T=F^{-1}(K_C)$.
For an effective divisor $D=\sum_zm_z z$ on $Z$ with
$\pi_*D\in Q$, define
\[
 E_{\max}(D)=\sum_{x\in C}
                 \max_{z\in\pi^{-1}(x)}m_z\,x.
\]
The dual of $dF_D$ is restriction of the pulled-back forms, so
\[
 \ker(dF_D)^*
       =H^0(C,K_C-E_{\max}(D)).
 \tag{F1T7.4}
\]
Here $E_{\max}\leq\pi_*D$ and
$\deg E_{\max}\geq\lceil4/3\rceil=2$.
The canonical section defining $\pi_*D$ supplies a nonzero
element of (F1T7.4). On a nonhyperelliptic genus-three curve,
a line bundle of degree at most two has at most one section,
unless it gives a degree-two pencil and hence a hyperelliptic
map. Consequently (F1T7.4) has dimension exactly one, and
$dF_D$ has rank two.

The finite map $\pi^{(4)}$ is flat: it is a finite map between
smooth varieties of the same dimension, so the local
Cohen--Macaulay criterion gives flatness. Its inverse image
of $Q$ has pure dimension two. The tangent dimension just
computed is two at every point; thus this scheme $T$ is smooth.
Let $p:T\to Q$. The kernel line in (F1T7.4) is precisely the
line of the canonical section defining $p(D)$. It therefore
gives the exact bundle sequence
\[
 0\longrightarrow p^*\mathcal O_Q(-1)
  \longrightarrow H^0(C,K_C)\otimes\mathcal O_T
  \longrightarrow N^\vee_{T/\operatorname{Sym}^4Z}
  \longrightarrow0.
 \tag{F1T7.5}
\]
This is the actual conormal morphism from the differential of
$F$, with no inference from dimensions alone.

Let $\eta$ be the class of the divisor $X_z$ of effective
degree-four divisors containing a fixed point $z\in Z$, and let
$\vartheta=a^*[\Theta_Z]$. The canonical class of the symmetric
power is $K_{\operatorname{Sym}^4Z}=2\eta+\vartheta$.
For an explicit normalization check, its tangent bundle is
the pushforward of $\mathcal O_{\mathcal D_Z}(\mathcal D_Z)$
for the universal divisor. The K-class is
$[\mathcal O(\mathcal D_Z)]-[\mathcal O]$. GRR on $Z$ gives
\[
 \operatorname{ch}(T_{\operatorname{Sym}^4Z})
           =6+(-2-\vartheta)e^\eta,
\]
using the mixed Poincare identity
$p_*(\gamma_Z^2)=-2\vartheta$. Its first character component
is $-2\eta-\vartheta$, proving the asserted canonical class.
The analogous point class restricts on $Q$ to
$c_1(\mathcal O_Q(1))$ and its pullback under $\pi^{(4)}$ is
$3\eta$. Taking determinants in (F1T7.5) and using adjunction
therefore gives the exact numerical identity
\[
                         K_T=(\vartheta-\eta)|_T.
 \tag{F1T7.6}
\]

\paragraph{2. Monodromy gives exactly three components.}
Use the canonical plane embedding $C\subset\mathbb P^2$ and
let $U\subset Q$ be the open set of lines with four distinct
intersection points on $C$. Consider the incidence surface
\[
 I_Z=\{(z,\ell):\pi(z)\in\ell\}
                 \longrightarrow Q.
\]
Projection to $Z$ makes $I_Z$ a $\mathbb P^1$-bundle, so it is
connected. Its restriction over $U$ is a connected etale
degree-twelve cover. Hence its monodromy group acts
transitively on the twelve points over a base line.

Label the four points on that line by $1,\ldots,4$ and their
three lifts by $\mathbb F_3$. Deck transformation translates
the three labels in every block. Monodromy commutes with it,
so the monodromy group lies in
$\mathbb F_3^4\rtimes\mathfrak S_4$.
At a generic point of the dual plane curve, a line has one
ordinary tangency. Its inertia exchanges two of the four
base points, and exchanges the corresponding three pairs
of lifts. It is therefore a colored transposition
$\rho_{ij}(a)$: the shifts in blocks $i,j$ are $a,-a$, and
the other blocks are fixed. Its square is the identity.

These inertia elements generate the monodromy group.
Indeed $Q=\mathbb P^2$ is simply connected, so every loop
in $U$ can be filled by a disc in $Q$; a small perturbation
meets the smooth part of the discriminant transversely and
avoids its finite singular locus. Cutting this disc expresses
its loop as a product of conjugate meridians.

Here is the complete resulting group calculation. The
graph on $\{1,2,3,4\}$ formed by the projected generating
transpositions is connected, since their action is transitive.
Choose a spanning tree and change the origins of the four
$\mathbb F_3$ labels to make its edge colors zero. These
tree transpositions generate the ordinary subgroup
$\mathfrak S_4$. There must be a generator with a nonzero
remaining color: otherwise the action on the twelve labels
would have the three separate sheet-label orbits, contradicting
transitivity. Multiplying such a $\rho_{ij}(a)$ by its ordinary
transposition produces the translation
$a(e_i-e_j)$, with $a\ne0$. Its $\mathfrak S_4$ conjugates
generate all of
\[
 V_0=\{(a_1,a_2,a_3,a_4)\in\mathbb F_3^4:
                         a_1+a_2+a_3+a_4=0\}.
\]
Every inertia element already lies in
$V_0\rtimes\mathfrak S_4$, so the monodromy group is exactly
\[
                         V_0\rtimes\mathfrak S_4=G(3,3,4).
 \tag{F1T7.7}
\]

The fibre of $T|_U\to U$ consists of the $3^4=81$ choices
of one lift of each of the four intersection points. Its
monodromy action is the action just computed on
$\mathbb F_3^4$. Its orbits are exactly the three values
of the coordinate sum; each orbit has size $3^3=27$.
Thus $T|_U$ has exactly three connected components.
Since $T$ is smooth and finite over $Q$, each of its
connected components is irreducible and dominates $Q$.
Removing the inverse image of the discriminant does not
disconnect it. Hence $T$ itself has exactly three components,
each of degree $27$.
The diagonal deck action adds $1$ to all four labels; it
adds $4=1$ in $\mathbb F_3$ to their sum. It therefore
permutes these three surfaces cyclically.

They map into the three distinct connected norm fibres.
Indeed $\pi_*H_1(Z,\mathbb Z)$ is the kernel of the
surjection $H_1(C,\mathbb Z)\to\mathbb Z/3$ defining the
cover. Its index is three, so the analytic norm kernel
has three connected components. The operator $1-\sigma$
is an isogeny on its identity component $P$, since
$H_1(P,\mathbb Q)$ has no invariant summand.
If $\sigma$ preserved a degree-four norm torsor, its
affine action on that torsor would have a fixed point:
the equation for such a point is solved by the
surjective isogeny $1-\sigma$ on $P$.
A $\sigma$-invariant line bundle on $Z$ admits a cyclic
linearization. Explicitly, rescale any isomorphism with
its $\sigma$ pullback by a cube root so that the product
of its three conjugates is the identity.
Etale descent then makes its degree divisible by three,
which excludes degree four. Thus $\sigma$ cyclically
permutes the three norm torsors. The three connected
$T_j$ map into them bijectively at the level of
connected-component labels.

\paragraph{3. Complete Abel fibres and their finite exceptional set.}
First, $Z$ is nonhyperelliptic. Otherwise its unique
hyperelliptic pencil is preserved by $\sigma$. The induced
automorphism of its target $\mathbb P^1$ has order three,
since the automorphisms over that target have order at most
two. Over a fixed point of this order-three automorphism,
the hyperelliptic fibre would be an invariant effective
divisor of degree two. Every orbit of the free action on $Z$
has length three, which makes that degree impossible.

It follows that every line bundle $L$ of degree four on $Z$
has $h^0(L)\leq2$. One can see this directly: three independent
sections, after removing the base divisor, give a nondegenerate
plane map of degree at most four. A birational plane image
would have arithmetic genus at most three. A map of degree
at least two must be a double cover of a conic, since the
nondegenerate plane image has degree at least two. This
would make $Z$ hyperelliptic.

The scheme-theoretic Abel fibre over $L$ is the complete
linear system $|L|=\mathbb P H^0(Z,L)$. If it meets $T$,
then $\operatorname{Nm}_\pi L=K_C$. For every $D\in|L|$,
the effective divisor $\pi_*D$ consequently lies in
$|K_C|$. Thus the whole Abel fibre is contained in $T$.
It is connected, so lies in one $T_j$. Every positive
dimensional fibre is therefore an actual smooth curve
\[
                         E=|L|\simeq\mathbb P^1.
 \tag{F1T7.8}
\]
The point-divisor bundle has
$\mathcal O(X_z)|_{|L|}\simeq\mathcal O_{\mathbb P H^0(L)}(1)$;
this follows from the universal divisor on
$Z\times\mathbb P H^0(L)$. Thus
$\eta E=1$, whereas $\vartheta E=0$.
Equation (F1T7.6) and adjunction give
\[
                         K_T E=-1,\qquad E^2=-1.
 \tag{F1T7.9}
\]

These fibres are finite in number. Indeed
$\mathcal O(X_z)$ is ample on $\operatorname{Sym}^4Z$:
its pullback to $Z^4$ is the external tensor product
of four degree-one ample line bundles, and ampleness
descends under a finite surjective map.
All curves (F1T7.8) therefore belong to the one
projective finite-type Hilbert scheme with polynomial
$r+1$ for this ample bundle (or its fixed Veronese
polynomial after taking a very ample power).
Their Hilbert tangent spaces are
\[
 H^0(E,N_{E/T_j})=H^0(\mathbb P^1,\mathcal O(-1))=0.
\]
They are isolated reduced points of that Hilbert scheme.
A finite-type scheme has only finitely many
zero-dimensional irreducible components, proving finiteness.
This argument concerns exactly the norm-compatible pencils;
it does not require a classification of all pencils on $Z$.

\paragraph{4. Blowdown and the descended embedding.}
The finitely many curves (F1T7.8) are disjoint, being distinct
fibres of the Abel map. Contract them by Castelnuovo's
contraction theorem, obtaining a smooth projective surface
\[
                         b:T_j\longrightarrow \overline S_j.
\]
An exact scheme-theoretic source is Stacks Project,
Lemma 54.16.8, tag \texttt{0C2M}; its definition of
contraction makes the target point regular of dimension two
and the source its blowup. The descended morphism below is
also the universal property in Lemma 54.16.1,
tag \texttt{0C5J}.
The Abel morphism is constant on every exceptional curve
and descends to a morphism
$g:\overline S_j\to\operatorname{Pic}^4Z$.
For example, near a contracted point choose an affine
neighborhood of its Abel value. Properness of $b$ gives
a neighborhood whose full inverse image maps into that
affine neighborhood; the coordinate functions descend
by $b_*\mathcal O_{T_j}=\mathcal O_{\overline S_j}$.
This proves the claimed factorization by an actual map.

The fibre description already shows that $g$ is injective
on points. We verify the differential also at every
contracted point. At $D\in|L|$, the divisor exact sequence
\[
0\longrightarrow\mathcal O_Z\longrightarrow L
 \longrightarrow L|_D\longrightarrow0
\]
identifies the tangent map of $a$ with its connecting map
to $H^1(Z,\mathcal O_Z)$. Its kernel is
$H^0(Z,L)/\mathbb C$. Along (F1T7.8) that is precisely
$T_DE$, of dimension one. The restriction of $da$ to
$T_DT_j$ therefore has rank one at every point of $E$.

Let $s=b(E)$. The normal derivative of a blowup is the
tautological inclusion
\[
 \mathcal O_E(-1)\longrightarrow
                     T_s\overline S_j\otimes\mathcal O_E.
 \tag{F1T7.10}
\]
In coordinates $y=xu$ it sends the normal vector
$\partial_x$ at $x=0$ to $(1,u)$; the other chart gives
the same tautological map. Since $a=g\circ b$, the
normal derivative of $a$ is the composite of (F1T7.10)
with $dg_s$. The rank-one calculation says that this
composite is nonzero on every fibre of $\mathcal O_E(-1)$.
Every line in $T_s\overline S_j$ occurs in (F1T7.10).
Thus $\ker dg_s=0$. Away from the exceptional curves,
$h^0(L)=1$ and the same divisor sequence gives injective
differential directly.

Consequently $g$ is a proper injective immersion, hence
a closed embedding. Its image is the reduced Abel
image $S_j$. This identifies $\overline S_j$ with
$S_j$ and proves the smooth-embedding assertion.
The actual map $T_j\to S_j$ is the indicated blowup.
Since its ordinary Abel fibres are single points,
its generic degree is exactly one.

\paragraph{5. Exact Abel degree and numerical record.}
Here we give the intersection calculation, also independently
checked by the root and fixed-surface contributions.
On $\operatorname{Sym}^4C$ write $x$ for the point class
and $\theta$ for the pullback of $\Theta_C$. The canonical
plane has class
\[
                         [Q]=x^2-x\theta+\theta^2/2.
 \tag{F1T7.11}
\]
Here is its direct degeneracy calculation. Let $\mathcal D$
be the universal divisor on
$C\times\operatorname{Sym}^4C$, with projections $q_C,q_S$,
and put
$\mathcal E=(q_S)_*(q_C^*K_C\otimes\mathcal O_{\mathcal D})$,
a rank-four bundle. The evaluation map
$H^0(K_C)\otimes\mathcal O\to\mathcal E$ has rank at most
two exactly on $Q$, because its kernel is $H^0(K_C-D)$.
This locus has the expected codimension two. Its
Thom--Porteous class is therefore $c_2(\mathcal E)$.
To compute it, use the universal-divisor K-class
$[q_C^*K_C](1-[\mathcal O(-\mathcal D)])$.
GRR on the curve factor gives
\[
 \operatorname{ch}(\mathcal E)
       =2-(-2-\theta)e^{-x},\qquad
 c_1(\mathcal E)=\theta-2x,\qquad
 \operatorname{ch}_2(\mathcal E)=x^2-x\theta.
\]
For the first equality, write
$c_1(\mathcal D)=4[\mathrm{pt}]+x+\gamma$;
the mixed Poincare component satisfies
$(q_S)_*(\gamma^2)=-2\theta$, as follows by expansion in
a symplectic basis. Integration of
$e^{4[\mathrm{pt}]}(1-e^{-c_1(\mathcal D)})
 (1-2[\mathrm{pt}])$ gives the displayed character.
The identity
$c_2=(c_1^2-2\operatorname{ch}_2)/2$ now gives (F1T7.11).
Let
$\Xi=a^*\operatorname{Nm}_\pi^*\Theta_C$ on
$\operatorname{Sym}^4Z$. Pullback of (F1T7.11) gives
\[
                         [T]=9\eta^2-3\eta\Xi+\Xi^2/2.
 \tag{F1T7.12}
\]
Relative to $\Theta_Z$, the endomorphism for
$\operatorname{Nm}_\pi^*\Theta_C$ is
$\pi^*\operatorname{Nm}_\pi=1+\sigma+\sigma^2$.
Its eigenvalues are $3$ on the dimension-three
invariant abelian factor and $0$ on the Prym.
The exact symmetric-product push formula and the
mixed theta intersections consequently give
\[
\begin{split}
 \int_{\operatorname{Sym}^4Z}\eta^2\vartheta^2&=42,\\
 \int_{\operatorname{Sym}^4Z}\eta\vartheta^2\Xi&=270,\\
 \int_{\operatorname{Sym}^4Z}\vartheta^2\Xi^2&=1080.
\end{split}
 \tag{F1T7.13}
\]
For the middle identity use
$\frac1{4!}\int_{J(Z)}\Theta_Z^6\operatorname{Nm}_\pi^*\Theta_C
=6!\cdot9/4!$; for the last use
$\frac1{3!}\int_{J(Z)}\Theta_Z^5
(\operatorname{Nm}_\pi^*\Theta_C)^2
=5!\cdot2!\cdot27/3!$.
Thus
\[
             \int_T\vartheta^2=378-810+540=108.
 \tag{F1T7.14}
\]
The deck action permutes the three $T_j$ and preserves
the Jacobian polarization, so each has value $36$.
Its Abel map has generic degree one by paragraph4,
proving (F1T7.3).

The accompanying normal-bundle Chern calculation gives
$K_{T_j}^2=15$, $c_2(T_j)=9$ and
$\chi(\mathcal O_{T_j})=2$. If $r_j$ is the number of
exceptional fibres that actually occur, the blowup map
proved above gives the exact identities
\[
 K_{S_j}^2=15+r_j,\qquad c_2(S_j)=9-r_j,\qquad
                 \chi(\mathcal O_{S_j})=2.
 \tag{F1T7.15}
\]
These formulas record $r_j$ as the actual finite count;
no asserted value of that count is used in the proof.
The deck action identifies the three counts.

\paragraph{6. The four requested frameworks and the actual maps.}
The FGA221 Hilbert scheme used here is precisely
$\operatorname{Hilb}^{r+1}(T_j,\mathcal O(X_z))$, with a
fixed very ample Veronese as required by that construction.
Its tangent at a pencil is the computed $H^0(N_{E/T_j})$.
For the final embedding $i_j:S_j\hookrightarrow P$, its
deformation map is the actual normal-sequence morphism
$H^1(P,T_P)\to H^1(S_j,N_{S_j/P})$, identified with the
dual semiregularity map as in F01--F02 and S01.
Thus the contraction construction supplies the precise
embedding to which that map is applied.

Over the nonhyperelliptic Hurwitz base, the construction of
$(Z,C,T,a)$ is algebraic. After a finite base change labels
the three components, its proper Abel push gives the
relative degree36 Chow family. On the nonempty flat open
of its scheme image it gives the actual Hilbert family of
the smooth embedded $S_j$. The equality
$a_*[T_j]=[S_j]$ has multiplicity one by paragraph4.
Its cycle class is horizontal in Betti and algebraic de Rham
realizations, with Gauss--Manin transport. This is the
specific variational realization of Grothendieck1966,
p.103, footnote13, through the FGA221 Hilbert/Chow maps.

For the SGA7 biextension maps, write
$n=\operatorname{Nm}_\pi:J(Z)\to J(C)$ and
$j=\pi^*:J(C)\to J(Z)$.
With the canonical principal polarizations one has
\[
 \widehat n\,\lambda_C=\lambda_Zj,\qquad
 \lambda_{n^*\Theta_C}
       =\widehat n\,\lambda_C n
       =\lambda_Z(1+\sigma+\sigma^2).
 \tag{F1T7.16}
\]
For $\iota:P\hookrightarrow J(Z)$ the induced polarization
is $\lambda_\beta=\widehat\iota\,\lambda_Z\iota$.
The normalized Poincare bundles satisfy the actual
pullback identity
$(n\times1)^*\mathcal P_C=(1\times\widehat n)^*\mathcal P_Z$
on $J(Z)\times\widehat{J(C)}$.
These are SGA7 I, VII2.9.5--2.9.6 and VIII3.2 in
the conventions of F03, and they prove the norm
polarization operator used in (F1T7.13).

Finally $Rb_*\mathcal O_{T_j}=\mathcal O_{S_j}$ for the
actual blowup, so
$Ra_*\mathcal O_{T_j}=i_{j,*}\mathcal O_{S_j}$ after
identifying its norm torsor with $P$ by translation.
GRR gives the full identity
\[
 \operatorname{ch}(i_{j,*}\mathcal O_{S_j})
       =i_{j,*}\bigl(\operatorname{td}(N_{S_j/P})^{-1}\bigr).
 \tag{F1T7.17}
\]
Its first nonzero component is
$\operatorname{ch}_2=[S_j]$.
SGA6 V, Appendix7.1--7.5 multiplies this component
by $m^2$ under $\psi^m$, with the higher components
given by the same displayed GRR formula.
The universal divisors, evaluation bundles, proper maps
and perfect K-classes above are finitely presented and
spread to a finitely generated ring after a finite
localization that includes inversion of three;
their Betti, de Rham and etale cycle realizations are
the concrete instances of F04.

The contraction and descent sources just cited were read
with their complete displayed proofs, including
Stacks tag \texttt{0C2L} for the regular contracted point.
The symmetric-product Chern normalizations were derived
above from the universal divisor. The intersections in
(F1T7.13) can be computed directly by pulling the forms
to the ordered product $Z^4$ and dividing by $4!$;
the complete mixed generating formula is
HD-F1-S09, (F1S9.9). Thus the required intersection
calculation does not depend on an additional preprint
theorem. The carried Poincare push formula for Abel images
also has the published locator Schoen (1988),
Corollary3.1 proof, printed p.25, citing
Griffiths--Harris, p.350.
No classification theorem for tetragonal curves is needed.
